\chapter{Clostridioides difficile Colitis}
\index{Clostridioides difficile Colitis}
\label{sec:Clostridioides difficile Colitis}

\section{Overview}
Clostridioides difficile colitis is an infectious inflammation of the colon caused by overgrowth of toxin-producing Clostridioides difficile bacteria, most commonly following antibiotic use. It is the leading cause of antibiotic-associated diarrhea and healthcare-associated infectious diarrhea, with approximately 500,000 infections annually in the United States.
\section{Classification}

\begin{description}[font=\normalfont\bfseries,leftmargin=3em,labelwidth=2.5em]
  \item[Cause] Infectious
  \item[Organ System] Gastrointestinal
  \item[Pathophysiology] Disruption of normal gut flora by antibiotics -> C. difficile overgrowth -> toxin A and B production -> colonic epithelial damage and pseudomembrane formation
  \item[Duration] Acute
  \item[Onset] Sudden
  \item[Mode of Transmission] Fecal-oral (spore-mediated)
  \item[Clinical Course] Variable; may recur
  \item[Severity] Mild to Severe
  \item[Extent of Disease] Localized to colon
  \item[Age Group] Adults (especially elderly)
  \item[Epidemiology] Highest incidence in hospitalized elderly patients; nosocomial spread
  \item[ICD-11 Code] 1A04
\end{description}

\section{Causes}
C. difficile colitis occurs when antibiotic therapy disrupts the normal colonic microbiome, allowing C. difficile spores to germinate and produce toxins A (enterotoxin) and B (cytotoxin), which cause direct colonic epithelial damage, inflammation, and formation of characteristic pseudomembranes. Clindamycin, fluoroquinolones, cephalosporins, and broad-spectrum penicillins are the most common antibiotic culprits. This condition develops from a combination of genetic and environmental factors that disrupt normal body processes. When these normal processes are disrupted, it leads to the symptoms and signs of the condition. Several factors can work together to trigger or make the condition worse. Knowing the specific cause helps your doctor choose the most effective treatment.

\section{Risk Factors}

\begin{itemize}[noitemsep]
  \item Recent or current antibiotic use
  \item Advanced age (over 65 years)
  \item Hospitalization or long-term care facility residence
  \item Immunocompromised state
  \item Proton pump inhibitor use
\end{itemize}

\section{Signs and Symptoms}

\subsection{Early symptoms}
\begin{itemize}[noitemsep]
  \item Early manifestations of clostridioides difficile colitis may be subtle and nonspecific
\end{itemize}

\subsection{Common symptoms}
\begin{itemize}[noitemsep]
  \item \subsection{Early symptoms}
  \item \subsection{Common symptoms}
  \item \textbf{Common symptoms:}
  \item Watery diarrhea occurring 3--15 times daily, often with a characteristic foul odor
  \item Abdominal cramping and diffuse lower abdominal tenderness
  \item Fever, nausea, and malaise
  \item Leukocytosis with left shift on laboratory studies
  \item Pain or discomfort in the affected area
  \item Difficulty performing daily activities
\end{itemize}

\subsection{Less common symptoms}
\begin{itemize}[noitemsep]
  \item Atypical or less common presentations of clostridioides difficile colitis may occur
\end{itemize}

\subsection{Emergency warning signs}
\begin{itemize}[noitemsep]
  \item Severe or worsening symptoms of clostridioides difficile colitis require immediate medical evaluation
\end{itemize}
\section{Progression and Stages}
Disease severity ranges from mild self-limited diarrhea to severe fulminant colitis. Severe disease is marked by elevated white blood cell count (>15,000 cells/mcL) and/or elevated serum creatinine (>1.5 times baseline), while fulminant colitis features toxic megacolon, colonic perforation, shock, or peritonitis requiring intensive care.

\section{Complications}

\begin{itemize}[noitemsep]
  \item Toxic megacolon with risk of colonic perforation
  \item Fulminant colitis requiring colectomy
  \item Recurrent C. difficile infection (20--30\% after first episode)
  \item Reduced quality of life
  \item Emotional and psychological effects
\end{itemize}

\section{Diagnosis}

\begin{itemize}[noitemsep]
  \item Stool testing for C. difficile toxins (toxin A/B enzyme immunoassay) or glutamate dehydrogenase (GDH) antigen
  \item Nucleic acid amplification test (NAAT) via PCR for toxin genes
  \item Lab tests to check how your organs are working
  \item Imaging tests
\end{itemize}

\section{Prevention}

\begin{itemize}[noitemsep]
  \item Judicious antibiotic stewardship to minimize unnecessary use
  \item Strict hand hygiene with soap and water (alcohol-based sanitizers do not kill spores)
  \item Go for regular check-ups so any problems are caught early
  \item Follow your doctor's screening recommendations
\end{itemize}

\section{Treatment Options}

\begin{itemize}[noitemsep]
  \item Vancomycin orally (125 mg four times daily) or fidaxomicin (200 mg twice daily) for initial episodes
  \item Metronidazole reserved for mild disease when vancomycin and fidaxomicin are unavailable
  \item Fecal microbiota transplantation for multiply recurrent C. difficile infection
  \item Lifestyle changes and self-care
  \item Surgery when needed (colectomy for fulminant disease)
\end{itemize}

\section{Living with It}

\begin{itemize}[noitemsep]
  \item Follow your treatment plan as prescribed
  \item Keep up with regular appointments with your healthcare provider
  \item Adjust your daily habits as needed
  \item Reach out to family, friends, or support groups for help
  \item Keep learning about your condition and how to manage it
\end{itemize}

\section{Outlook}
Most patients with C. difficile colitis recover completely with appropriate antibiotic therapy, though recurrence rates remain significant. Your outlook depends on the specific condition, how advanced it is when found, and how well it responds to treatment. Getting treated early generally leads to better results. Many people manage this condition effectively with the right treatment. Regular check-ups are important to monitor your condition and adjust treatment as needed.
